\section{How much can one seed tell you?}
\label{sec:seeds}

Everything in the report for exercises 3.1 and 3.2 came from one seed per experiment (31 for the robot, 32 for the student). This is common in course projects, and to be honest also in a part of the literature \cite{henderson2018}. Here we rerun the exact same experiments, same functions and same order of random draws, on 100 seeds (seeds 1 to 100, so the two seeds of the report are among them). We checked first that the seeds of the report give back the numbers of the report to the last digit, otherwise the comparison would not mean much.

We look at it from three angles. First we go through the report and pick the sentences that make a claim we can check on one run (``$\eta = 0.01$ gives the smallest error'', ``the three values of $\alpha$ end on the optimal policy'', and so on), and count on how many of the 100 seeds the sentence would still be true. Second, for a few comparisons between two settings, we ask how many seeds one has to average before the ranking agrees with the 100-seed ranking most of the time (Colas et al.\ \cite{colas2018} do a real power analysis for this kind of question, we just bootstrap from our 100 seeds), and we report the probability of improvement $P(X < Y)$ between one random run of each setting, as suggested by Agarwal et al.\ \cite{agarwal2021}. Third, we look at the shape of some distributions across seeds, because a mean $\pm$ standard deviation can be misleading when the distribution is skewed or has two modes.
